\section*{Chapitre \ref{ch:b2:equilibrium-shifts} --- Constantes d'équilibre et déplacements}

\begin{solution}{exo:b2:equilibrium-shifts:1}
$K^\circ = \exp(32\,800/(8.314 \times 298.15)) = \num{5.6e5}$ et, pour
$+\qty{32.8}{kJ/mol}$, son inverse, \num{1.8e-6}. Un facteur 100 exige
$\Delta(\Delta_r G^\circ) = -RT\ln 100 = \qty{-11.4}{kJ/mol}$.
\end{solution}

\begin{solution}{exo:b2:equilibrium-shifts:2}
$\ln K^\circ(400) = \ln(5.4 \times 10^5) + (91\,880/8.314)(1/400 -
1/298.15) = 13.20 - 9.44 = 3.76$, $K^\circ \approx 43$. Les tables donnent 35.6~:
sur \qty{100}{K}, l'approximation de $\Delta_r H^\circ$ constante est déjà
en erreur de 20\,\% (la réaction devient plus \omterm{def:b2:reaction-enthalpy:exothermic}{exothermique} quand $T$ augmente).
\end{solution}

\begin{solution}{exo:b2:equilibrium-shifts:3}
(a) $n = 1$, $r = 0$, $\varphi = 2$~: $v = 1$ (la pression de vapeur est une fonction
de $T$). (b) $v = 3 - 1 + 2 - 3 = 1$. (c) $v = 3 - 1 + 2 - 1 = 3$. (d) Quatre
espèces, une réaction, deux phases (le carbone et le gaz)~: $v = 4 - 1 + 2 - 2 =
3$~; si le gaz provient uniquement du carbone et de la vapeur d'eau, $n(\ce{CO}) = n(\ce{H2})$
dans le gaz, $s = 1$ et $v = 2$.
\end{solution}

\begin{solution}{exo:b2:equilibrium-shifts:4}
Chauffage~: sens \omterm{def:b2:reaction-enthalpy:exothermic}{endothermique}, vers \ce{NO2} (sens direct). Compression~:
vers moins de molécules de gaz, \ce{N2O4} (sens inverse). Argon à volume constant~:
aucun effet, les pressions partielles des gaz qui réagissent sont inchangées.
\end{solution}

\begin{solution}{exo:b2:equilibrium-shifts:5}
À \qty{298}{K}~: $\Delta_r G^\circ = 180.6 - 298.15 \times 0.0248 =
\qty{173.2}{kJ/mol}$, $K^\circ = \num{4.5e-31}$. À \qty{2000}{K}~: $\Delta_r
G^\circ = 180.6 - 2000 \times 0.0248 = \qty{131.0}{kJ/mol}$, $K^\circ =
\num{3.8e-4}$. Avec $\Delta\nu_{\text{gaz}} = 0$, $x(\ce{NO})^2 = K^\circ x(\ce{N2})
x(\ce{O2})$~: $x(\ce{NO}) = \sqrt{3.8 \times 10^{-4} \times 0.78 \times 0.21} =
\num{7.9e-3}$, presque un pour cent. Les moteurs brûlent à de telles températures~; quand les
gaz refroidissent, $K^\circ$ chute énormément, mais la décomposition de \ce{NO} est
très lente~: l'équilibre chaud est figé dans les gaz d'échappement (jusqu'à ce qu'un catalyseur
l'élimine).
\end{solution}

\begin{solution}{exo:b2:equilibrium-shifts:6}
Avec $\xi$ mol d'ester et une quantité totale constante, $Q =
\xi^2/[(1 - \xi)(n_B - \xi)]$. Pour $n_B = 1.0$~: $\xi^2/(1 - \xi)^2 = 4$, $\xi
= \qty{0.67}{mol}$. Pour $n_B = 3.0$~: $3\xi^2 - 16\xi + 12 = 0$, $\xi =
\qty{0.90}{mol}$. Éliminer l'eau maintient $Q$ au-dessous de $K^\circ$ quel que soit $\xi$~: la
réaction continue d'avancer jusqu'à consommation de l'acide.
\end{solution}

\begin{solution}{exo:b2:equilibrium-shifts:7}
Sans argon~: quantités $1 - \alpha$, $2\alpha$, total $1 + \alpha$~;
$Q = 4\alpha^2/(1 - \alpha^2) = 0.148$, $\alpha = \sqrt{0.148/4.148} = 0.189$.
Avec \qty{1}{mol} d'argon sous \qty{1}{bar}~: total $2 + \alpha$, $Q =
4\alpha^2/[(1 - \alpha)(2 + \alpha)] = 0.148$, donc $4.148\alpha^2 + 0.148\alpha
- 0.296 = 0$, $\alpha = 0.250$~: diluer à pression totale constante abaisse les
pressions partielles, ce qui favorise le côté qui compte le plus de molécules de gaz. À
volume constant, les pressions partielles des gaz qui réagissent ne changent pas~:
$\alpha$ reste égal à 0.189.
\end{solution}

\begin{solution}{exo:b2:equilibrium-shifts:8}
Solide seul~: $n = 3$, $r = 1$, $s = 1$ ($p(\ce{NH3}) = p(\ce{HCl})$),
$\varphi = 2$~: $v = 3 - 1 - 1 + 2 - 2 = 1$~; on choisit $T$, et la pression est
fixée (une «~pression de dissociation~»). Avec de l'ammoniac ajouté, $s = 0$ et $v = 2$~:
on peut choisir $T$ et une pression~; le produit $p(\ce{NH3})\,p(\ce{HCl})$ est
fixé par $T$.
\end{solution}

\begin{solution}{exo:b2:equilibrium-shifts:9}
$\Delta_r H^\circ = -R\ln(K_2/K_1)/(1/T_2 - 1/T_1) = -8.314 \times
\ln(0.0173)/(1/600 - 1/500) = \qty{-101}{kJ/mol}$, plus négative qu'à
\qty{298}{K} (\qty{-91.9}{kJ/mol}), comme le prévoyait la \omterm{thm:b2:reaction-enthalpy:kirchhoff}{loi de Kirchhoff}
($\Delta_r C_p^\circ < 0$).
\end{solution}

\begin{solution}{exo:b2:equilibrium-shifts:10}
$Q = \dfrac{n_{\ce{NH3}}^2\,n_{\text{tot}}^2}{n_{\ce{N2}}n_{\ce{H2}}^3}
\left(\dfrac{p^\circ}{p}\right)^2$. À $T$, $p$ constantes~:
$\partial\ln Q/\partial n_{\ce{N2}} = -1/n_{\ce{N2}} + 2/n_{\text{tot}}$, positive
lorsque $n_{\ce{N2}} > n_{\text{tot}}/2$, c'est-à-dire $x(\ce{N2}) > 1/2$. Alors $Q$
dépasse $K^\circ$ et l'équilibre se déplace dans le sens inverse. À volume constant,
$\partial\ln Q/\partial n_{\ce{N2}} = -1/n_{\ce{N2}} < 0$~: toujours dans le sens direct.
\end{solution}

\begin{solution}{exo:b2:equilibrium-shifts:11}
$K^\circ = \exp(-1620/(8.314 \times 1100)) = 0.84$~: $p(\ce{CO2}) =
\qty{0.84}{bar}$ à l'équilibre, soit $n = pV/RT = 0.84 \times 10^5 \times
0.0100/(8.314 \times 1100) = \qty{0.092}{mol}$ de gaz. Avec \qty{0.050}{mol}
de calcaire, tout se décompose ($p = \qty{0.46}{bar} < K^\circ p^\circ$)~:
pas d'équilibre, il y a rupture d'équilibre. Avec \qty{0.200}{mol}, équilibre~:
\qty{0.092}{mol} de \ce{CO2} et de \ce{CaO}, et il reste \qty{0.108}{mol} de \ce{CaCO3}.
\end{solution}

\begin{solution}{exo:b2:equilibrium-shifts:12}
Le lit 1 sort vers \qty{890}{K} avec 68\,\% de conversion, le lit 2 vers
\qty{785}{K} avec 92\,\%, le lit 3 vers \qty{725}{K} avec 98\,\%. Dans un seul
lit adiabatique, le gaz s'échauffe en réagissant jusqu'à rencontrer la courbe
d'équilibre, qui à cette température ne permet qu'environ 68\,\%. Un catalyseur froid
tout du long est exclu parce que le catalyseur n'agit qu'à chaud~; refroidir
entre les lits préserve la vitesse et relève la limite d'équilibre. Le dernier lit
est limité par l'équilibre à la plus basse température à laquelle le catalyseur
agit encore (les usines absorbent alors le trioxyde et font passer le gaz une fois de plus sur
un catalyseur).
\end{solution}

\begin{solution}{pb:b2:equilibrium-shifts:1}
\textbf{1.} $\Delta_r H^\circ = \qty{-91.88}{kJ/mol}$, $\Delta_r S^\circ =
2(192.77) - 191.61 - 3(130.68) = \qty{-198.1}{J/(K.mol)}$, $\Delta_r G^\circ =
\qty{-32.8}{kJ/mol}$, $K^\circ = \num{5.6e5}$.
\textbf{2.} $K^\circ(700) = \exp(-54\,380/(8.314 \times 700)) = \num{8.8e-5}$~;
$K^\circ(800) = \exp(-77\,320/(8.314 \times 800)) = \num{8.9e-6}$.
\textbf{3.} $\Delta_r H^\circ = -R\ln(K_{800}/K_{700})/(1/800 - 1/700)$,
soit \qty{-106}{kJ/mol}, plus négative qu'à \qty{298}{K}.
\textbf{4.} $\ln K^\circ(723) = \ln(8.75 \times 10^{-5}) + 12\,777 \times
(1/723.15 - 1/700)$, avec $12\,777 = 106\,230/8.314$~: $K^\circ =
\num{4.9e-5}$.
\textbf{5.} $\Delta_r G^\circ(723) = -91.88 + 723.15 \times 0.1981 =
\qty{51.4}{kJ/mol}$, $K^\circ = \num{1.9e-4}$~: quatre fois trop grande, car
$\Delta_r H^\circ$ et $\Delta_r S^\circ$ varient sur \qty{425}{K}
($\Delta_r C_p^\circ \approx \qty{-44}{J/(K.mol)}$).
\textbf{6.} \ce{N2} $1 - \xi$, \ce{H2} $3 - 3\xi$, \ce{NH3} $2\xi$, total $4 -
2\xi$.
\textbf{7.} $x = 2\xi/(4 - 2\xi)$, donc $1 - x = (4 - 4\xi)/(4 - 2\xi)$~; la
fraction de diazote $(1 - \xi)/(4 - 2\xi) = (1 - x)/4$, et celle de dihydrogène
est trois fois plus grande.
\textbf{8.}
\[
  K^\circ = \frac{x^2}{\frac{1-x}{4}\left(\frac{3(1-x)}{4}\right)^3}
  \left(\frac{p^\circ}{p}\right)^2 = \frac{256}{27}\,\frac{x^2}{(1 - x)^4}
  \left(\frac{p^\circ}{p}\right)^2 ;
\]
il reste à prendre la racine carrée.
\textbf{9.} $a = (\sqrt{27}/16) \times \sqrt{4.9 \times 10^{-5}} \times 1 =
\num{2.27e-3}$~; $x/(1 - x)^2 = a$ donne $x = 0.0023$.
\textbf{10.} $a = 0.455$~; $ax^2 - (2a + 1)x + a = 0$~: $x = 0.25$.
\textbf{11.} $n = 3$, $r = 1$, $s = 1$ (rapport 1 : 3 conservé par la réaction),
$\varphi = 1$~: $v = 2$. Une fois $T$ et $p$ choisies, plus rien n'est libre.
\textbf{12.} Pression~: sens direct ($\Delta\nu_{\text{gaz}} = -2$). Température~:
sens inverse (réaction \omterm{def:b2:reaction-enthalpy:exothermic}{exothermique}).
\textbf{13.} Un gaz inerte à pression totale constante abaisse les pressions
partielles des gaz qui réagissent, comme le ferait une baisse de pression~: sens inverse.
\textbf{14.} $\partial\ln Q/\partial n_{\ce{N2}} = (-1/0.30 + 2)/n_{\text{tot}} < 0$~:
$Q$ passe sous $K^\circ$, sens direct.
\textbf{15.} Non~: à volume constant, $\partial\ln Q/\partial n_{\ce{N2}} =
-1/n_{\ce{N2}} < 0$ quelle que soit la composition~: sens direct.
\textbf{16.} Presque sans ammoniac, $Q$ est très inférieur à $K^\circ$ à l'entrée~:
le gaz entre loin de l'équilibre et réagit dans le sens direct.
\textbf{17.} Le gaz passe trop peu de temps sur le catalyseur pour que
l'équilibre soit atteint~; un contact plus long exigerait un convertisseur plus grand et plus
cher pour un faible gain.
\textbf{18.} $x = 2\xi/(4 - 2\xi) = 0.18$ donne $\xi = \qty{0.305}{mol}$ par
mole de \ce{N2} introduite~: \omterm{def:b2:equilibrium-shifts:single-pass}{conversion par passage} du dihydrogène (et du diazote)
$0.305$, soit 31\,\%.
\textbf{19.} En régime permanent, toute l'alimentation fraîche finit par être convertie~:
\qty{2.00}{mol} d'ammoniac sortent par unité de temps~; le convertisseur doit traiter $1.00/
0.305 = \qty{3.3}{mol}$ de \ce{N2} par unité de temps, le reste étant recyclé.
\textbf{20.} L'argon (et le méthane) qui entrent avec l'alimentation ne réagissent jamais et
ne sont pas extraits avec l'ammoniac liquide~: sans \omterm{def:b2:equilibrium-shifts:single-pass}{purge}, ils
s'accumuleraient dans la boucle et abaisseraient les pressions partielles des réactifs.
\textbf{21.} À \qty{500}{K}, le catalyseur au fer est trop lent~: l'équilibre
serait favorable, mais il ne serait jamais approché en un temps industriel.
\textbf{22.} À \qty{723}{K} et sous \qty{200}{bar}, alimentation stœchiométrique, gaz
parfaits~: $\boldsymbol{x(\ce{NH3}) \approx 0.25}$.
\end{solution}
